\chapter[Estudio de base: metodología]{Estudio de base: datos y metodología}
\label{cap:metodologia}

Este capítulo documenta los experimentos realizados en el proyecto \emph{vision}
hasta el 24 de septiembre de 2026. Constituyen el estudio de base para
caracterizar TA1C, delimitar espacios de soluciones y establecer referencias
de calidad y costo. El código,
las configuraciones, las predicciones y las métricas están en
\path{metodologia/}. La guía de reproducción indica qué recursos externos
se necesitan.

\section{Tareas, corpus y particiones}

En detección se clasifica una publicación que promociona una noticia. La
entrada principal es el texto de la publicación, o \emph{teaser}, y la salida
indica si hay clickbait. En spoiling se usa también el artículo para generar
una respuesta breve a la pregunta que la publicación deja abierta. Las tareas
se evalúan por separado, porque detectar bien el clickbait no garantiza una
respuesta fiel a la noticia.

Los experimentos principales usan los archivos oficiales de TA1C para
IberLEF 2025 \citep{ta1c2025}. Detección cuenta con 2.800 ejemplos de
entrenamiento, 700 de desarrollo y 700 de prueba; spoiling, con 300, 100 y
100. Los CSV se convierten a JSONL sin perder los identificadores, las
etiquetas, las particiones ni las entradas. También se guardan las cantidades
y las huellas SHA-256 de los archivos originales para verificar su origen.
Una exploración anterior usó otra versión de detección, de
2.100/700/700 ejemplos. Sus cifras se presentan por separado.

Los modelos textuales habituales se entrenan en \emph{train} y se comparan
en \emph{dev}, que tiene 700 ejemplos. Los experimentos visuales D10 y D11
dividen \emph{train+dev} (3.500 ejemplos) en cinco grupos estratificados.
Cada ejemplo se predice con un modelo entrenado en los otros cuatro grupos;
estas son las predicciones fuera de fold u OOF. Sus puntuaciones se comparan
entre sí, no con las obtenidas solo en desarrollo. El conjunto de prueba
oficial contiene solo teaser y metadatos: no trae el artículo ni la imagen.
Se reserva para una evaluación textual final, aunque existen ejecuciones
históricas de líneas base sobre él que se señalan en los resultados. Una
evaluación final multimodal necesita un conjunto independiente con las
modalidades disponibles, fijado antes de seleccionar modelos y separado por
noticia y posibles duplicados.

\section{Modalidades y diseños experimentales}

Las líneas base D0 combinan TF-IDF con clasificadores lineales. D1 usa
\emph{embeddings} sin ajustarlos; D2--D4 ajustan codificadores de texto de
distintos tamaños; D5--D6 prueban clasificadores basados en modelos de
lenguaje grandes. Para spoiling, S0 usa reglas sencillas; S1 extrae texto;
S3 usa modelos generativos abiertos; y S6, servicios por API. En
\path{metodologia/experiments/pareto/} se registran la configuración,
las semillas, el tiempo, los recursos y las predicciones de cada corrida.
Las demás variantes del catálogo aún no se han ejecutado.

Los experimentos visuales comparan una variante con texto (T) y otra con
texto e imagen (TI) de la misma familia de modelos. Se usa una imagen por
ejemplo: primero la de la noticia y, si falta, la del tuit. Cuando no hay
imagen, se registra su ausencia y se usa solo texto. D10a prueba
representaciones visuales sin ajustarlas; D10b combina RoBERTa-BNE y SigLIP2
con una red MLP residual. D11a usa una MLP para fusionar las entradas y D11b
emplea atención cruzada entre palabras y fragmentos de imagen. En D11,
el codificador visual permanece fijo. El régimen \emph{fixed} mantiene la
referencia textual mientras entrena la fusión; \emph{joint} ajusta también
el codificador textual con una pérdida auxiliar. Ambas arquitecturas se
probaron con una semilla y cinco grupos. En spoiling, S7 compara un mismo
generador con y sin imagen o pie de foto. D12 está preparado para combinar
publicación, artículo e imagen en un espacio común, pero aún no tiene una
comparación completa.

No todas las imágenes estaban disponibles o eran pertinentes. El registro
incluye archivos faltantes o pequeños y posibles duplicados entre particiones.
Los grupos de D10 y D11 mantienen la proporción de clases, pero no separan
las noticias o imágenes duplicadas. La comparación mide el cambio observado
al añadir imágenes en estas condiciones; no permite atribuirlo por sí solo
al contenido visual ni generalizarlo a otros medios.

\section{Métricas, incertidumbre y costo}

En estos experimentos de detección se prioriza la F1 de la clase clickbait,
que corresponde a la puntuación publicada en el tablero oficial. También se
guardan precisión, exhaustividad, F1 macro, PR-AUC y la matriz de confusión;
la F1 macro es una medida complementaria. Para ordenar las corridas de spoiling se usa
una implementación interna de ROUGE-L basada en la subsecuencia común más
larga de tokens Unicode en minúsculas. La implementación local de BERTScore F1
con \texttt{bert-base-multilingual-cased}, sin reescalado, fue la única que
se aproximó a la cifra de la línea base oficial. Los valores locales de
ROUGE-L y BLEU no se comparan con la tabla oficial de 2025: su programa de
evaluación no está disponible y las variantes probadas no reprodujeron
esas cifras.

Las diferencias entre T y TI se calculan sobre los mismos identificadores,
con la misma etiqueta y, si corresponde, el mismo grupo de validación. Se
informa la diferencia TI menos T, tanto para todos los ejemplos como para
los que tienen imagen. El intervalo percentil del 95\,\% se obtiene con
2.000 remuestreos pareados y una semilla fija. Si incluye cero, los datos
no respaldan una mejora concluyente. El intervalo refleja la variación entre
ejemplos, pero no toda la que proviene de volver a entrenar, elegir modelos
o recuperar otras imágenes.

El frente de Pareto compara calidad en desarrollo y operaciones de cómputo
estimadas por ejemplo (FLOPs). Una configuración domina a otra si cuesta
igual o menos y logra una calidad igual o mayor, con al menos una ventaja
estricta. Detección y spoiling se grafican por separado y solo incluyen
corridas comparables dentro de cada tarea: 700 ejemplos de detección o 100
de spoiling, con FLOPs conocidos. Los
servicios por API se muestran aparte, con su costo en dólares por 1.000
ejemplos. Los FLOPs no miden el tiempo de respuesta, que depende del equipo
y de sus condiciones de uso. El frente describe las corridas disponibles,
pero no basta para elegir un sistema de producción.

\section{Reproducibilidad y límites del estudio de base}

La carpeta \path{metodologia/} reúne el código, las configuraciones, las
entradas normalizadas, el registro de corridas, las predicciones y el
proyecto Power BI. Las figuras se generan con
\path{metodologia/scripts/graficar_resultados.py}. Para repetir D11 con los
mismos píxeles hace falta el archivo de imágenes usado en Colab, que no está
en la copia local. Se localizó una copia candidata en Drive, todavía sin
cotejar su hash y cobertura con las corridas. Descargar de nuevo las imágenes
desde sus URL podría producir entradas distintas. Los siguientes pasos son revisar duplicados entre grupos, repetir
la fusión con más semillas, fijar una evaluación independiente que disponga
de imágenes para los modelos multimodales y pedir a personas que juzguen la
fidelidad y utilidad de los spoilers. La prueba oficial de detección solo
permite evaluar modelos que usen sus entradas textuales.
